\scanpage{77}
Alors, définissant $T_0$ par (8) (pour $t\in\Sigma_{J}^\ast$ donné) on a un isom canonique
\[
T_0^{\otimes2}\simeq T_{\Sigma,t}\tag*{(12)}
\]

\deleted{\unsure{une fois|ceci posé|en fait} \illegible{2} \unsure{supposé donné un} \unsure{$Q_0$} \unsure{$J$|de} \illegible{1} \unsure{qui soit|quand}} \inserted{Ceci posé, on appellera \uline{épinglage de Legendre}} de $E$, subordonné au carré $Q_0$ (muni des données (11)) \deleted{\unsure{tout}} \inserted{un} isomorphisme
\[
T=T_{E,0}\xleftarrow{\overset{k}{\sim}}T_0 \tag*{(13)}
\]
compatible avec les hom. norme, i.e. rendant commutatif le diagramme
\diagram{\begin{tikzcd}[ampersand replacement=\&,column sep=small,row sep=large]
\text{(14)} \& T_0^{\otimes2}\arrow[rr,"{\overset{k^{\otimes2}}{\sim}}"] \arrow[dr,"{(12)}"] \arrow[d,phantom,"{\unsure{!}}"] \& \& T^{\otimes2}\arrow[dl,"{(5)}"]\\
\& {} \& T_{\Sigma,t} \& {}
\end{tikzcd}}

Il y en exactement deux, différant par le signe i.e. formant un torseur sous $\mu_2$.

Bien entendu, on veut pouvoir définir la notion d'épinglage de $E$ (en termes d'un $Q_0$ donné une fois pour \unsure{toutes|toute}) pour $E$ variable -- mais alors $\underline{\unsure{\omega|w}}=\underline\omega_J$ n'est pas encore déterminé (puisque $J$ ne l'est pas). Donc la donnée d'un épinglage de Legendre sera défini comme le couple \inserted{$(\alpha,k)$} formé d'un isom.
\[
\alpha:\underline\omega_J\xleftarrow{\sim}D_{Q_0}=Q_0/\{\pm1\} \tag*{(15)}
\]
\uline{et} d'un isom. $k$ (13) satisfaisant la compatibilité (14),
% archive sheet number: 76
